\subsection{Voting Scheme Defense}

The defense mechanism employs $n \geq 3$ parallel processing units, each
independently computing the tracking error $e(k)$.
This approach follows the redundant voting paradigm used in
safety-critical systems~\cite{lyons1962tmr, siewiorek1998reliable}.
Under nominal conditions, all units compute the error identically.
An adversarial attack corrupts at most one unit by injecting additive noise
into its error computation. This redundancy enables detection and exclusion
of the corrupted processing unit through pairwise comparison.

\centering{WORK IN PROGRESS}
\par
Consider $n = 3$ processing units computing the tracking error. Unit $i$ produces
\begin{equation}
	\label{eq:channel_error}
	e_i(k) = e(k) + \delta_i(k),
\end{equation}
where $\delta_i(k)$ represents the noise on processing unit $i$.
Clean units are ideal, with $\delta_i(k) = 0$, which gives
\begin{displaymath}
	e_i(k) = e(k)\text{.}
\end{displaymath}


% The pairwise disagreement metric between units $i$ and $j$ is defined as
% \begin{equation}
% 	\label{eq:pairwise_disagreement}
% 	V_{ij}(k) = \|e_i(k) - e_j(k)\|^2.
% \end{equation}
%
% The minimum-variance pair is selected by
% \begin{equation}
% 	\label{eq:pair_selection}
% 	(i^*, j^*) = \underset{(i,j)}{\arg\min} \; V_{ij}(k).
% \end{equation}
%
% The defended error signal is the average of the selected pair,
% \begin{equation}
% 	\label{eq:defended_error}
% 	\hat{e}(k) = \frac{e_{i^*}(k) + e_{j^*}(k)}{2}.
% \end{equation}
%
% Substituting $\hat{e}(k)$ into the control law from \eqref{eq:control_function}
% gives the defended control input
% \begin{equation}
% 	\label{eq:defended_control}
% 	\hat{u}(k) = g(x(k))^{-1} \bigl( -f(x(k)) + \Delta r(k) - \beta \operatorname{sgn}(\hat{e}(k)) \bigr),
% \end{equation}
% and the defended closed-loop state equation
% \begin{equation}
% 	\label{eq:defended_closed_loop}
% 	x(k+1) = x(k) - \beta \operatorname{sgn}(\hat{e}(k)) + \Delta r(k).
% \end{equation}
